\chapter[Band 38: Monoide]{Band 38 auf einen Blick}
\noindent\textbf{Monoide}\par\medskip
Dieser Band ergänzt Halbgruppen um ein neutrales Element.
Die Aussagen zu Teilbarkeit verwenden Kommutativität; bei
\glqq prim impliziert irreduzibel\grqq{} kommt Kürzbarkeit hinzu.
Potenzprodukte und Produktquadrate stammen aus Band~28.

\begin{BandUebersicht}
\textbf{Monoidaxiome}
\UebFormel{\begin{aligned}
&\MonoidAlg A\star e\ \leftrightarrow\\
&\SemigroupAlg A\star\land e\in A\\
&\quad\land\forall a\in A\ (e\star a=a=a\star e).
\end{aligned}}
\UebQuellen{\FormulaRefAuto{MonoidDef}[def]}
&
\textbf{Eindeutigkeit und Beispiele}
\UebFormel{\begin{aligned}
&\MonoidAlg A\star e,\ \MonoidAlg A\star u
 \ \vdash e=u,\\
&\CommutativeMonoidAlg{\mathbb N}{+}{0},\\
&\CommutativeMonoidAlg{\mathbb N}{\cdot}{1}.
\end{aligned}}
\UebQuellen{\FormulaRefAuto{MonoidIdentityUnique}[thm];
\FormulaRefAuto{NaturalAdditiveCommutativeMonoid}[thm];
\FormulaRefAuto{NaturalMultiplicativeCommutativeMonoid}[thm]} \\
\addlinespace[.8ex]

\textbf{Einheiten}
Für ein Monoid \((A,\star,e)\):\UebFormel{\begin{aligned}
&\MonoidUnit{A,\star,e}u\ \leftrightarrow\\
&u\in A\land\exists v\in A\
 (u\star v=e=v\star u).
\end{aligned}}
\UebQuellen{\FormulaRefAuto{MonoidUnitDef}[def]}
&
\textbf{Einheiten der Potenzhalbgruppe}
Die große Potenzhalbgruppe ist ein Monoid mit Einheit \(\{e\}\).
Ihre invertierbaren Elemente sind genau die Einermengen invertierbarer
Grundelemente:\UebFormel{\begin{aligned}
&X\in\PowNE A\ \vdash\\
&\MonoidUnit{\PowNE A,\star_{\mathcal P},\{e\}}X\\
&\quad\leftrightarrow\exists u\in A\
 (X=\{u\}\land\MonoidUnit{A,\star,e}u).
\end{aligned}}
\UebQuellen{\FormulaRefAuto{PowerSemigroupMonoid}[thm];
\FormulaRefAuto{PowerMonoidUnitCharacterization}[thm]} \\
\addlinespace[.8ex]

\textbf{Teilbarkeit im kommutativen Monoid}
Bei \(\CommutativeMonoidAlg A\star e\) und \(a,b\in A\):\UebFormel{\begin{aligned}
&a\mid b\ \leftrightarrow\exists c\in A\ (b=a\star c).
\end{aligned}}
\UebQuellen{\FormulaRefAuto{CommutativeMonoidDividesDef}[def]}
&
\textbf{Einheiten als Teiler der Eins}
\UebFormel{\begin{aligned}
&\CommutativeMonoidAlg A\star e,\ u\in A\\
&\vdash\MonoidUnit{A,\star,e}u
 \ \leftrightarrow u\mid e.
\end{aligned}}
\UebQuellen{\FormulaRefAuto{CommutativeMonoidUnitIffDividesIdentity}[thm]} \\
\addlinespace[.8ex]

\textbf{Prim und irreduzibel}
Im kommutativen Monoid ist \(p\) irreduzibel, wenn \(p\) keine
Einheit ist und jede Zerlegung \(p=a\star b\) einen Einheitsfaktor besitzt.
Es ist prim, wenn es keine Einheit ist und\UebFormel{\begin{aligned}
&a,b\in A,\ p\mid a\star b\\
&\qquad\vdash p\mid a\lor p\mid b.
\end{aligned}}
\UebQuellen{\FormulaRefAuto{CommutativeMonoidIrreducibleDef}[def];
\FormulaRefAuto{CommutativeMonoidPrimeDef}[def]}
&
\textbf{Prim impliziert irreduzibel}
\UebFormel{\begin{aligned}
&\CommutativeMonoidAlg A\star e,\\
&\CancellativeSemigroupAlg A\star,\\
&\MonoidPrime{A,\star,e}p\\
&\qquad\vdash\MonoidIrreducible{A,\star,e}p.
\end{aligned}}
\UebQuellen{\FormulaRefAuto{CommutativeCancellativeMonoidPrimeIsIrreducible}[thm]} \\
\addlinespace[.8ex]

\textbf{Monoidmorphismen}
Für Monoide \(\mathfrak A=(A,\star,e_A)\),
\(\mathfrak B=(B,\diamond,e_B)\) bedeutet
\(\MonoidHom f{\mathfrak A}{\mathfrak B}\):\UebFormel{\begin{aligned}
&f:A\to B,\quad f(e_A)=e_B,\\
&x,y\in A\ \vdash\
 f(x\star y)=f(x)\diamond f(y).
\end{aligned}}
\UebQuellen{\FormulaRefAuto{MonoidHomDef}[def];
\FormulaRefAuto{MonoidIsoDef}[def]}
&
\textbf{Komposition und Einheiten}
Die Komposition zweier Monoidmorphismen ist wieder ein
Monoidmorphismus. Ein Halbgruppenisomorphismus zwischen Monoiden
erhält automatisch die neutralen Elemente und die Einheiten:\UebFormel{\begin{aligned}
&f(e_A)=e_B,\\
&\MonoidUnit{A,\star,e_A}u\\
&\qquad\leftrightarrow\MonoidUnit{B,\diamond,e_B}{f(u)}
 \quad(u\in A).
\end{aligned}}
\UebQuellen{\FormulaRefAuto{MonoidHomComposition}[thm];
\FormulaRefAuto{SemigroupIsoBetweenMonoidsIdentityAndUnits}[thm]} \\
\addlinespace[.8ex]

\textbf{Monoidales Produktquadrat}
Für eine nichtleere Halbgruppe \((A,\star)\) sei
\(A^2=\{a\star b\mid a,b\in A\}\).
Wenn \(\MonoidAlg{A^2}\star e\), setzt man\UebFormel{\begin{aligned}
&r_e:A\to A^2,\qquad r_e(a)=e\star a.
\end{aligned}}
\UebQuellen{\FormulaRefAuto{SemigroupSquareDef}[def];
\FormulaRefAuto{MonoidSquareInflationDef}[def]}
&
\textbf{Rückführung auf das Quadrat}
Unter diesen Voraussetzungen gilt:\UebFormel{\begin{aligned}
&r_e|_{A^2}=\operatorname{id}_{A^2},\\
&a,b\in A\ \vdash\
 a\star b=r_e(a)\star r_e(b).
\end{aligned}}Ferner besitzt \((\PowNE A)^2\) genau dann ein neutrales Element,
wenn \(A^2\) eines besitzt.
\UebQuellen{\FormulaRefAuto{MonoidSquareInflationProperties}[thm];
\FormulaRefAuto{PowerSemigroupSquareMonoidCriterion}[thm]} \\
\addlinespace[.8ex]

\textbf{Fasern der Rückführung}
Für die vorige Rückführung und \(g\in A^2\):\UebFormel{\begin{aligned}
&C_g=\{a\in A\mid r_e(a)=g\}.
\end{aligned}}
\UebQuellen{\FormulaRefAuto{MonoidSquareInflationDef}[def]}
&
\textbf{Potenzfasern}
\UebFormel{\begin{aligned}
&\{X\in\PowNE A\mid
 \{e\}\star_{\mathcal P}X=\{g\}\}\\
&\qquad=\PowNE{C_g}.
\end{aligned}}Die Faser über einer Einermenge ist somit genau die
nichtleere Potenzmenge der entsprechenden Grundfaser.
\UebQuellen{\FormulaRefAuto{MonoidSquarePowerFiber}[thm]} \\
\addlinespace[.8ex]
\end{BandUebersicht}

